\documentclass[12pt]{article}

% Fluxo do campo para um cone
% Mexilhoeira Grande, Sex Nov 30 20:19:52 WET 2018
% Written by Tordar

\usepackage{graphicx}
\usepackage{pst-all}
\usepackage{pst-3dplot}



\begin{document}

\pagestyle{empty}

\begin{pspicture}
\rput(2,0){
\psset{linewidth=0.75mm,linecolor=black,linestyle=dashed,dash=4mm 1mm,arrows=->}
\pstThreeDLine[arrows=-](0,0,0)(-6,0,0)
\pstThreeDLine[arrows=-](0,-6,0)(0,0,0)
\pstThreeDLine[linestyle=solid](0,0,0)(6,0,0)
\pstThreeDLine[linestyle=solid](0,0,0)(0,6,0)
\pstThreeDLine[linestyle=solid](0,0,3)(0,0,5)
\pstThreeDTriangle[linewidth=1mm,linecolor=green,fillstyle=solid,fillcolor=yellow,arrows=-](0,-2,0)(0,0,4)(0,2,0)
\pstThreeDLine[linewidth=1.5mm,linecolor=blue,linestyle=solid,arrows=->](-5,0,0)(-3,0,0)
\pstPlanePut[plane=xz](6.7,0,0){\scalebox{1.5}{$X$}}
\pstPlanePut[plane=yz](0,6.1,0){\scalebox{1.5}{$Y$}}
\pstPlanePut[plane=yz](0,0,5.5){\scalebox{1.5}{$Z$}}
\pstThreeDPut[plane=yz](-4,0,1){\scalebox{1.5}{$\mathbf{E}$}}
%\pstPlanePut[plane=yz](0,-0.2,1.5){\scalebox{1.5}{$A$}}
}
\end{pspicture}

\end{document}
